\chapter{Preface}
仅为一家之言, 仓促完成且未经他人审校, 难免存在疏漏, 谬误. 恳请各位读者不吝指正, 共同探讨与完善. 

本笔记存档于 \href{https://github.com/IceySwan/Notes}{GitHub}, 如发现任何错误请提 \href{https://github.com/IceySwan/Notes/issues/new}{issue} 或联系 \href{mailto:hi@icey.one}{hi@icey.one}.

\section*{Contents}
本笔记旨在记录笔者对 Riemann-Hilbert 方法的理解, RH 方法作为一种强大的数学工具, 在非线性偏微分方程的求解, 尤其是可积系统中具有重要的应用价值. 通过对其理论背景, 基本方法及具体应用的学习, 笔者希望为读者提供一个清晰的入门指引. 各章节主要内容如下: 

第一章主要介绍 RH 方法的背景知识, 包括 Plemelj 定理, RH 问题等基础知识. 对 RH 方法已有初步了解的读者可跳过本章. 

第二章主要介绍利用 RH 方法求解零边界的 NLS 方程, 通过构造特征函数, 分析其解析性与对称性, 以及建立相关的 RH 问题，最终推导出 NLS 方程的 N 孤子解. 本章节主要参考了复旦大学范恩贵老师的讲义, 并感谢同门韩刻蓉的帮助与讲解. 

第三章主要介绍利用 RH 方法求解反时间, 反空间, 及反时空的反演形式, 这些方程是耦合 NLS 方程在不同约束条件下的特殊形式. 通过 Lax 对出发, 分析其散射数据的对称性，进一步推导出这些非局部方程的 $ N $ 孤子解. 跳过了传统 RH 方法寻找 Jost 函数和跳跃矩阵的过程, 直接利用约化条件对散射数据进行分析, 极大简化了推导过程. 

第四章为在第三章的基础上, 将结论进一步推广到三维的反时空 NLS 方程. 通过分析其散射数据的对称性，推导出三维反时空 NLS 方程的 N 孤子解. 相比于第三章, 第四章的结论可以推广至 $N \times N$ 的情形, 是本笔记的精髓. 本章节构成了我的第一篇论文. 感谢君言师姐推导的约化, 助益很多. 

\section*{感想}
转瞬已是毕业季. 思索再三, 决定将我硕士期间学习 Riemann-Hilbert 方法的笔记整理成册, 以飨读者. 回想当初动笔时还是在研一在老家过年期间, 今天在同一张桌子上完成笔记初稿, 感慨颇多. 2024 年春季, 我第一次来到郑州大学, 还不知将与可积系统结缘. 经过了一学期的学习, 我对 RH 方法有了初步的了解, 便萌生了将所学内容整理成册的想法, 便开始了这段漫长的写作旅程. 

整个写作过程并不算顺利, 期间经历了许多挑战和困难. 首先是数学上的挑战, RH 方法作为一种强大的数学工具, 理论艰深复杂, 每每提笔写作方知基础尚有欠缺, 第二章作为最先动笔的章节, 历经多次修改已不见初稿样貌. 其次是时间上的挑战, 作为学生, 还需兼顾学业和科研, 写作只能利用课余时间和假期, 进度缓慢. 在此奉劝各位有志于科研的同仁, 不要担任各类职位, 俗务缠身实难找出完整时间科研 . 其三是生活上的挑战, 适逢春节期间, 老屋年久失修, 外墙坍塌和两间屋顶损毁, 白日和父亲砌墙, 应付家庭琐事; 唯有夜间缓慢推进写作. 农村老家没有暖气, 零下十余度的环境实难称舒适, 写作需要物理意义上的热情与字面意义上的争分夺秒. 深夜伏案, 四野寂寂, 唯余键盘声在黑暗中回荡. 恍若独行于漫长甬道, 徘徊在幽邃迷宫. 然正如 A. Zee 所言:「夜航人自有夜行法」. 而我也悟到了适合自己的写作方式: 白天构思提纲, 夜至十点后开机写作, 睡前再反思增补. 

笔者体悟出一种难以言喻的, 某种漠然的相互理解, 如越过高墙, 漫步于满月下的林地, 并在夜色中获得慰籍. 是邪，非邪？解释或属妄诞，感受毕竟真实. 

\begin{flushright}
    Icey Swan\\
    \zhtoday\\
    于河南老家深夜
\end{flushright}